\documentclass[../main.tex]{subfiles}
\begin{document}

\chapter*{Course Map}
\addcontentsline{toc}{chapter}{Course Map}
\markboth{Course Map}{}
\label{ch:coursemap}

\begin{figure}[H]
  \centering
  \includegraphics[width=\linewidth]{Overview.png}
  \caption{The course as two parallel tracks. Green nodes form the conceptual track, red nodes the empirical track, and blue-grey nodes the integrative topics that join them.}
  \label{fig:overview}
\end{figure}

Green (conceptual track): language of research → conceptual design (W2) → technical report (W6) — what is studied and how the argument is framed. Red (empirical track): research steps → sampling and data (W3) → validation (W4) → data management (W5) — how evidence is gathered, checked, and managed. Both tracks converge at \textbf{Data Management}, the node where the conceptual plan and the empirical discipline meet.

\end{document}
